\documentclass[12pt]{article}

\usepackage{geometry}
\usepackage{graphicx}
\usepackage{hyperref}

\geometry{margin=1in}

\title{News Credibility Analysis Using Logistic Regression}
\author{Advik Khandelwal}
\date{\today}

\begin{document}

\maketitle

\begin{abstract}
This project presents a machine learning based News Credibility Analysis
system that predicts whether a news article is real or fake. Natural
Language Processing (NLP) techniques and Logistic Regression are used
to classify textual news content. The system performs preprocessing,
feature extraction using TF-IDF, and classification with high accuracy.
\end{abstract}

\section{Introduction}

The rapid spread of misinformation through online platforms has made
automatic fake news detection an important research problem. Manual
verification of news articles is time-consuming, motivating the use
of machine learning techniques for automated credibility analysis.

This project builds a model that classifies news articles as real or fake
based on textual content.

\section{Dataset}

The dataset was obtained using KaggleHub and consists of news articles
containing the following fields:

\begin{itemize}
\item Title
\item Text content
\item Label (Real or Fake)
\end{itemize}

The title and article text were combined to form the final input content.
Missing values were removed before training.

\section{Data Preprocessing}

Text preprocessing was performed using Natural Language Processing
techniques:

\begin{itemize}
\item Conversion of text to lowercase
\item Removal of URLs and special characters
\item Retaining alphanumeric characters
\item Stopword removal using NLTK
\end{itemize}

A cleaned text column was generated and used for feature extraction.

\section{Feature Extraction}

TF-IDF (Term Frequency–Inverse Document Frequency) vectorization was used
to convert text into numerical features.

The configuration used:

\begin{itemize}
\item Maximum features: 20,000
\item N-gram range: (1,2) (unigrams and bigrams)
\item Minimum document frequency: 5
\item Maximum document frequency: 0.7
\item Sublinear term frequency scaling enabled
\end{itemize}

\section{Model Training}

The dataset was split into training and testing sets using an
80:20 ratio with stratified sampling.

A Logistic Regression classifier was trained with:

\begin{itemize}
\item Maximum iterations: 2000
\item Class weight: balanced
\item Regularization parameter (C): 2.0
\end{itemize}

Logistic Regression was chosen due to its efficiency and strong
performance in binary text classification problems.

\section{Model Evaluation}

Model performance was evaluated using precision, recall,
F1-score, and confusion matrix metrics.

\subsection{Classification Report}

\begin{verbatim}
              precision    recall  f1-score   support

           0       0.97      0.96      0.96      7006
           1       0.96      0.97      0.97      7302

    accuracy                           0.96     14308
   macro avg       0.96      0.96      0.96     14308
weighted avg       0.96      0.96      0.96     14308
\end{verbatim}

\subsection{Confusion Matrix}

\begin{verbatim}
[[6726  280]
 [ 232 7070]]
\end{verbatim}

The model achieves an overall accuracy of approximately 96\%, indicating
strong predictive capability and balanced performance across both classes.

\section{Model Deployment}

The trained Logistic Regression model and TF-IDF vectorizer were saved
using Python's pickle module. These saved models were later integrated
into a Streamlit web application that allows users to input news text
and receive real-time credibility predictions.

\section{Conclusion}

This project demonstrates the effectiveness of machine learning and
NLP techniques for fake news detection. Logistic Regression combined
with TF-IDF feature extraction provides accurate and computationally
efficient classification results.

\section{Future Work}

Future improvements may include:

\begin{itemize}
\item Using deep learning models such as LSTM or Transformers
\item Expanding dataset diversity
\item Deploying the system on cloud platforms
\item Adding explainable AI visualizations
\end{itemize}

\end{document}